\chapter{Background}
\label{chap:background}

% Drafted under T-018, 13 August 2026. Section source records, one per section, are in
% notes/chapter2-source-records/ and name the claim each citation supports.
% Section list is fixed by notes/dissertation-structure-v0.3.md (revision of 13 August 2026).
% OPEN, tracked as T-028: G1-G4 in 2.7 carry no \label. Chapter 1 section 1.3 does not yet
% point at them and Chapter 8 does not yet close them. Whether they become macros in
% dissertationmacros.tex, as the research questions did under D-002, or a \label/\ref pair,
% is the author's decision and is still open.
% Page budget under D-007, ceiling 14 pages:
%   2.1 3.0, 2.2 2.5, 2.3 2.0, 2.4 1.75, 2.5 1.25, 2.6 1.25, 2.7 1.5.
% D-008 fixes the formalism split. The incidence matrix and the vertex and hyperedge degree
% matrices are defined in 2.2. Section 2.3 carries no equations. The hypergraph Laplacian,
% the convolution operator and both attention equations belong to Chapter 5.
% Two figures, each about a third of a page and already inside the budget above: an APK
% artefacts and analysis modes figure in 2.1, and a pairwise graph against hypergraph lifting
% figure in 2.2 marking the clique-expansion collapse.
% Section 2.3 is written neutrally on attention and commits to no causal claim, so that it
% stands whichever way T-012 is decided.
% T-023 / D-004, 20 September 2026: shared performance measures added within Section 2.6.
% The seven main sections and the evaluation label are retained. Claim-to-source records
% and the inventory of reported versus proposed metrics are in the T-023 writing ticket.

This chapter sets out the material the later chapters share: the artefacts an Android
application exposes to analysis, the graph and hypergraph representations used to describe
programs, the learning operators defined over those representations, detection in networked and
cyber-physical systems, adversarial machine learning, and the methodology by which any of it is
assessed. Positioning against competing systems is not attempted here. Every contribution
chapter carries its own account of the work specific to its method, and that is where published
systems are compared.

Learning-based security analysis has a common shape. Samples are collected and labelled, a
representation is derived from each sample, a model is fitted, its output is thresholded, and
the result is measured on data held back from training. A methodological defect can enter at any
of these stages, and the defects that recur have been catalogued stage by stage
\cite{arp2022dos}. Four assumptions carried over from other application domains do not hold in
this setting. Training and test data are taken to come from one stationary distribution, which
the evolution of malware denies. The available sample is taken to represent the deployment,
which its composition and its timeline usually deny. Errors are treated as symmetric and the
classes as roughly balanced, whereas operational traffic is neither \cite{sommer2010outside,
axelsson2000base}.
The data are assumed not to be chosen by an opponent. Section~\ref{sec:background:evaluation}
takes up the first three and Section~\ref{sec:background:adversarial} the fourth.

The Android sections open the chapter, because Android malware analysis is the centre of the
dissertation and because a hyperedge is easier to read once function-call graphs have been
defined. One consequence of that ordering should be signalled here.
Section~\ref{sec:background:networks} supports Chapter~\ref{chap:anomaly} although it appears
after the Android material, so a reader following chapter order meets the foundations later than
the chapter that draws on them. Section~\ref{sec:anomaly:1} repeats the pointer.

\section{Android Applications and Their Analysis}
\label{sec:background:android}

Chapters~\ref{chap:hybroid} and~\ref{chap:hgann} both take an Android application package as
input and derive features from it. This section fixes what that package contains, what each mode
of analysis is able to observe in it, and the two mechanisms that most often prevent observation
altogether. Network-flow representation is deferred to
Section~\ref{sec:background:networks}, and the graph objects the two chapters build are defined
in Section~\ref{sec:background:graphs}.

\subsection{Dalvik Bytecode, Manifest, and Permissions}
\label{sec:background:android:artefacts}

An Android application is distributed as an application package, a ZIP archive whose
security-relevant members are a manifest and one or more files of compiled bytecode
\cite{liu2023deeplearning}. Source written in Java or Kotlin is compiled to bytecode for the
Dalvik virtual machine \cite{arp2014drebin} and stored in \texttt{classes.dex}, with further
numbered files such as \texttt{classes2.dex} for larger applications. The container is specified by the
platform \cite{aosp2026dexformat}, as is the instruction set \cite{aosp2026dalvikbytecode}.
The Android Runtime (ART) replaced Dalvik as the default runtime from Android~5 and executes
DEX bytecode \cite{aosp2026art}. Its instructions encode both opcodes and operands. Operand
fields can contain register indices or literal values and can reference other DEX items,
such as methods and strings \cite{aosp2026dalvikbytecode}. Section~\ref{sec:hybroid:5}
describes an embedding that retains the opcode while omitting this operand information.

The manifest declares the components an application provides, the permissions it requests, the
intents its components filter, and the hardware features it expects. It is recoverable without
any reasoning about code, which is what makes it the cheapest security-relevant metadata
available: four of Drebin's eight feature sets are taken from the manifest and the remaining
four from the disassembled bytecode \cite{arp2014drebin}.

Access to sensitive functionality is mediated by permissions, each carrying a protection level,
of which only the level marked \emph{dangerous} guards capability able to cause real harm
\cite{felt2011android, aosp2026permissions}. Two properties of this system bear on the later
chapters. The mapping from a permission to the API calls it guards is documented incompletely
and has to be recovered by analysing the platform: PScout extracted it from roughly 3.4 million
lines of Android source across four releases and found more than 75 permissions whose guarded
API the documentation did not describe \cite{au2012pscout}. The specification is also unstable,
since fine-grained least privilege trades against stability across releases
\cite{au2012pscout}, and deprecation on the platform side has been measured as one source of the
drift reported for permission-based detection \cite{sabbah2026understanding}.
Section~\ref{sec:hgann:5} relies on both properties, because it selects permissions by
protection level and maps flagged API calls back to their call sites.

\subsection{Static, Dynamic, and Hybrid Analysis}
\label{sec:background:android:analysis}

An application may be analysed without running it, by running it, or by combining the two.
Figure~\ref{fig:background:apk} sets the three modes against the artefacts of the previous
subsection. The static mode dominates practice: across the studies of deep learning for Android
malware defence reviewed by Liu et al., 73 per cent worked statically, 17 per cent dynamically,
and 10 per cent combined both \cite{liu2023deeplearning}.

\begin{figure}[tbp]
  \centering
  % Approved imagegen variant; trim removes only the temporary preview label.
  \includegraphics[width=\linewidth,trim=0 0 0 115bp,clip]
    {background-apk-analysis.png}
  \caption[Android package contents and analysis modes]{What an application package holds and
  what each mode of analysis reaches. Static observation covers all shipped code and not what
  the code does; dynamic observation covers what the code did, but only along the paths the
  harness triggered, and it alone reaches network traffic, which is the channel
  Chapter~\ref{chap:hybroid} uses. The graph objects built from the bytecode are named in
  Section~\ref{sec:background:graphs}.}
  \label{fig:background:apk}
\end{figure}

Static analysis reads the package and reasons about code that may never execute, which yields
coverage of everything shipped at modest computational cost. The engineering is mature.
Disassembly and call-graph construction are supplied by standard toolkits
\cite{desnos2012android, valleerai2010soot}, and taint analysis that models Android's component
lifecycle is available \cite{arzt2014flowdroid}. What no amount of engineering removes is a
ceiling on what static reasoning can decide. An adversary is able to construct obfuscation whose
analysis reduces to an NP-hard problem, which defeats semantics-aware static detection
\cite{moser2007limits}. That construction was demonstrated on x86 binaries, and the argument
carries to Dalvik by the same reasoning rather than having been shown there directly.

Dynamic analysis observes the application while it runs. Instrumentation ranges from taint
tracking on the device \cite{enck2014taintdroid} to whole-system emulation that reconstructs the
operating-system and Dalvik views together \cite{yan2012droidscope}. Traffic recorded during
execution is a further channel, and it is the one Chapter~\ref{chap:hybroid} uses. Two limits
apply throughout. Only the paths the harness triggers are observed, and an application may
detect the analysis environment and withhold the behaviour under investigation
\cite{petsas2014rage, vidas2014evading}.

Neither mode dominates the other. Measured under one protocol, static analysis proved at least
as effective as dynamic analysis, depending on how the applications were stimulated during
execution, while the combination of the two measured best \cite{onwuzurike2018family}. Hybrid
pipelines are consequently an established design \cite{lindorfer2015marvin}, and the question
Chapter~\ref{chap:hybroid} takes up is what such a combination buys on one set of samples under
a single protocol.

\subsection{Obfuscation, Packing, and Repackaging}
\label{sec:background:android:obfuscation}

Obfuscation is ordinary developer practice as much as attacker practice. Of 1.7 million free
applications on Google Play, 24.92 per cent had been obfuscated by their developers
\cite{wermke2018large}, and a measurement across 114,560 applications drawn from Google Play,
third-party markets and malware corpora found string encryption more frequent in malware and
packing more frequent outside Google Play \cite{dong2018understanding}. The effect on detection
is measurable. Three thousand benign and three thousand malicious applications, expanded into
73,362 obfuscated variants by 29 strategies, left most commercial anti-malware products severely
affected, and trivial transformations sufficed \cite{hammad2018large}. A later study of twelve
detectors found that none held its ideal-setting performance once obfuscation, concept drift and
adversarial examples were introduced \cite{gao2024comprehensive}.

Packing is a distinct mechanism and is kept separate here. A packed application ships an
encrypted original together with a wrapper that decrypts and loads it at runtime, so static
extraction fails outright instead of degrading. Commercial packers are used by malware authors,
and custom packers are written specifically to defeat the available unpackers
\cite{duan2018things}.

The packaging model also allows an attacker to graft a malicious rider onto a benign carrier
\cite{li2017understanding}. Applications produced this way are near-identical outside the
injected payload, and corpora assembled from them therefore contain duplicates. What those
duplicates cost depends on the task: the measured effect is limited for supervised
classification and substantially larger for unsupervised family clustering
\cite{zhao2021impact}. Section~\ref{sec:background:evaluation} returns to this as a problem of
evaluation rather than of extraction.

\section{Graph Representations of Programs}
\label{sec:background:graphs}

Representations are defined here, before the operators that learn over them. Two families are
needed. Chapters~\ref{chap:hybroid} and~\ref{chap:hgann} both derive a graph from program
structure, and Chapter~\ref{chap:hgann} then lifts that graph to a hypergraph. The notation
introduced in this section is the notation those chapters use, and neither redefines it.

\subsection{Control-Flow and Function-Call Graphs}
\label{sec:background:graphs:pairwise}

A \emph{control-flow graph} is a directed graph over the basic blocks of a single procedure. Its
vertices are maximal instruction sequences that contain no branch except possibly at the end,
and a directed edge records that control may pass from one block to the next. The object is
intra-procedural by definition \cite{bilot2024survey}. A \emph{call graph} takes the procedures
of a program as its vertices, with an edge from $f$ to $g$ recording that a call site in $f$ may
invoke $g$. It is inter-procedural, and a statically recovered call graph over-approximates,
since the analysis admits edges that no execution takes. For Android the call graph is taken at
method granularity over the recovered bytecode, and this dissertation calls that object a
\emph{function-call graph}: vertices are methods, and directed edges are method invocations
\cite{gascon2013structural, norouzian2025hgannmal}. The two objects are kept distinct throughout,
because their vertex sets and their edge semantics differ.

The two compose. Embeddings computed on per-method control-flow graphs may be attached to the
vertices of a function-call graph, which makes intra-procedural and inter-procedural structure
available together \cite{bilot2024survey}. Dropping the call edges from that composition leaves
a third object: a set of independently embedded control-flow graphs, pooled to application
level, with no invocation structure at all. The three are named here so that each contribution
chapter can state which of them it builds without introducing new machinery, and so that a
claim made about one is not read as a claim about another.

Whichever object is chosen, it is an artefact of the analyser rather than a property of the
program. Recovery depends on the toolkit and on how the component lifecycle is modelled
\cite{valleerai2010soot, arzt2014flowdroid, desnos2012android}, and abstracting call targets to
package or family level shrinks the vertex set at the cost of the callee's identity
\cite{onwuzurike2019mamadroid}. Vertex features may equally be learned rather than hand-built,
from the block structure of a procedure \cite{dai2016discriminative, xu2017neural} or from the
body of a method \cite{alon2019code2vec}.

\subsection{Hypergraphs and Higher-Order Relations}
\label{sec:background:graphs:hypergraphs}

A hypergraph $\mathcal{H} = (\mathcal{V}, \mathcal{E}, w)$ has a finite vertex set $\mathcal{V}$
and a family $\mathcal{E}$ of subsets of $\mathcal{V}$, each subset a hyperedge $e$ carrying a
weight $w(e) > 0$. Membership is recorded by the incidence matrix
$H \in \{0,1\}^{|\mathcal{V}| \times |\mathcal{E}|}$, where $h(v,e) = 1$ exactly when $v \in e$.
The degree of a vertex is $d(v) = \sum_{e \in \mathcal{E}} w(e)\, h(v,e)$ and the degree of a
hyperedge is $\delta(e) = \sum_{v \in \mathcal{V}} h(v,e) = |e|$. Writing $D_v$ and $D_e$ for the
diagonal matrices of vertex and hyperedge degrees and $W$ for the diagonal matrix of hyperedge
weights completes the apparatus \cite{zhou2006learning}. A graph is the special case in which
$\delta(e) = 2$ for every hyperedge, so every graph is a hypergraph and the converse fails.

An edge relates exactly two vertices, and encoding a group of several as a set of pairwise edges
loses the fact that its members belonged to one group. The projected graph cannot distinguish a
single hyperedge on three vertices from three separate pairwise relations among the same
vertices \cite{zhou2006learning}. Two projections are standard. \emph{Clique expansion} replaces
each hyperedge by the complete graph on its members, and \emph{star expansion} introduces a new
vertex per hyperedge joined to each member \cite{agarwal2006higherorder}.
Figure~\ref{fig:background:hypergraph} shows the hyperedge memberships and the information
lost in their pairwise projection.

\begin{figure}[tbp]
  \centering
  % Approved imagegen variant; trim removes only the temporary preview label.
  \includegraphics[width=\linewidth,trim=0 0 0 42bp,clip]
    {background-hypergraph-projection.png}
  \caption[Pairwise projection of a hypergraph]{The same six methods represented as
  (a) a hypergraph, (b) its incidence matrix $H$, and (c) its unweighted clique expansion.
  A matrix entry is one precisely when the method belongs to the corresponding hyperedge.
  Hyperedge $e_3$ joins two vertices and is therefore an ordinary graph edge. The lower
  example shows that the triangle over $m_1, m_2, m_3$ arises both from the single hyperedge
  $e_1$ and from three separate pairwise relations. The projection preserves pairwise
  adjacency but does not uniquely determine the original hyperedge memberships
  \cite{zhou2006learning}.}
  \label{fig:background:hypergraph}
\end{figure}

The sceptical reading of these definitions deserves stating, because it is the stronger position
to argue from. Agarwal et al.\ show that the hypergraph learning formulations of their period
reduce to the Laplacian of one expansion or the other, that the two coincide for uniform
hypergraphs, and that for non-uniform ones they differ only in how hyperedges of unequal size
are weighted; their conclusion is that graphs, not hypergraphs, lie at the heart of the problem
\cite{agarwal2006higherorder}. Chitra and Raphael sharpen the boundary. A random walk on a
hypergraph whose vertex weights do not vary between the hyperedges a vertex belongs to, the
\emph{edge-independent} case, is equivalent to a walk on a weighted clique expansion, so an
operator built on such a walk carries no higher-order information. Only \emph{edge-dependent}
vertex weights escape the equivalence \cite{chitra2019randomwalks}.

What the projection costs has since been measured rather than asserted. An evaluation framework
that expands, learns and then reconstructs reports depressed performance at vertex, hyperedge
and hypergraph level across eight datasets \cite{kang2024low}; two recurring hyperedge patterns
have been identified as the source of the loss in every projection, together with a
combinatorial impossibility result on recovering the discarded structure \cite{wang2024graphs};
and an expressivity hierarchy names the invariants a clique expansion cannot represent
\cite{jiang2026widthwall}. Section~\ref{sec:background:learning} describes the operators built on
this apparatus, and Chapter~\ref{chap:hgann} develops the one it uses.

\section{Learning over Graphs}
\label{sec:background:learning}

The operators described here learn representations over a structure whose construction was
fixed in Section~\ref{sec:background:graphs}. Three families are relevant: neighbourhood
aggregation on pairwise graphs, attention-weighted aggregation, and convolution over an
incidence structure. This section states what each does and what is known about its limits. The
notation stays as defined above, and the operator Chapter~\ref{chap:hgann} builds on is
developed there rather than here.

\subsection{Message Passing}
\label{sec:background:learning:messagepassing}

Most graph neural networks share one scheme. Each vertex holds a state, that state is updated
from the states of its neighbours by a permutation-invariant aggregation, and the update is
applied for a fixed number of rounds \cite{gilmer2017neural}. The convolutional instance fixes
the aggregation coefficients in advance from the degrees of the two endpoints, so a vertex
combines its neighbours in proportions that the graph determines and training does not alter
\cite{kipf2017semi}.

Two limits on this family are established. The first is a bound on what it can tell apart: no
network in the family distinguishes two graphs that one-dimensional Weisfeiler-Leman colour
refinement assigns the same colours, and an architecture reaching that bound is as
discriminative as the scheme permits \cite{xu2019powerful}. The second concerns depth. Repeated
propagation drives the representations within a connected component towards one another, so
adding rounds to widen a vertex's receptive field costs discriminability between vertices
\cite{li2018deeper}. Both bear on program graphs, where the structures compared are often large
and locally similar.

\subsection{Attention}
\label{sec:background:learning:attention}

Attention replaces fixed coefficients with learned ones. A scoring function is applied to each
vertex and its neighbour, the scores across a neighbourhood are normalised, and the resulting
weights determine the combination \cite{velickovic2018graph}; the mechanism was introduced for
sequences before it was carried to graphs \cite{vaswani2017attention}. The motivation is that
some neighbours matter more than others, and that which ones matter should be learned rather
than read off the degree.

Three results qualify that motivation, and they should be kept apart. Brody et al.\ show that
the ranking the standard formulation induces over a vertex's neighbours is the same for every
query vertex, so the mechanism selects globally important neighbours rather than neighbours
important to a particular vertex; a revised scoring function removes the restriction
\cite{brody2022gatv2}. Fang et al.\ bound both scoring functions, the linear form used by the
original mechanism and the multi-layer form introduced by its revision, showing that each
realises only finitely many neighbour rankings at bounded width and depth \cite{fang2025kaa}.
The same work reports that on node-property benchmarks over pairwise graphs, evaluated with
GCN, GraphSAGE and GIN as non-attentive baselines, no unmodified attentive model outperformed
all of those baselines on any dataset or task \cite{fang2025kaa}. The defensible reading of that
last result is negative about the evidence rather than about the mechanism: an advantage for
attention over fixed aggregation has not been demonstrated on those benchmarks, and the
expressive power of the standard scoring functions is bounded.

Whether either finding transfers to operators defined over hyperedges is not established in the
literature, since both analyses concern pairwise operators. Section~\ref{sec:hgann:8} shows that
the static-ranking argument carries over to the additive hyperedge score used there. For other
scoring functions over hyperedges, and for the finite-ranking bound, the question is open, and
Section~\ref{sec:background:learning:hypergraphconvolution} returns to it.

A methodological point follows, and it concerns how such architectures are designed rather than
how they perform. Where a model changes its relational encoding and its aggregation weighting at
the same time, the contribution of each change cannot be recovered from the reported totals. The
two have to be varied separately for either to be attributed.

\subsection{Hypergraph Convolution}
\label{sec:background:learning:hypergraphconvolution}

Operators over a hypergraph propagate in two stages. Signal moves from the vertices to the
hyperedges that contain them, is combined at each hyperedge, and returns to the member vertices,
with the incidence structure and the degree matrices of
Section~\ref{sec:background:graphs:hypergraphs} supplying the normalisation
\cite{feng2019hypergraph, gao2023hgnn}. Two design routes exist. One reduces the hypergraph to a
graph and applies a pairwise operator to the result \cite{yadati2019hypergcn}, which inherits
the losses recorded in Section~\ref{sec:background:graphs:hypergraphs}; the other treats a
hyperedge as a multiset and learns the functions that aggregate it
\cite{chien2022allset}. Attention has been carried across to this setting by making the
incidence structure itself learnable, so that membership is weighted by a real value instead of
a binary indicator \cite{bai2021hypergraphattention}.

Three observations bound what is known. Standard hypergraph convolution averages over the
members of a hyperedge and therefore attenuates variation between them, a low-pass bias whose
discarded high-frequency component has been characterised \cite{li2026highpass}. The operator
most widely used has since been superseded by a stability-oriented successor from its own
authors \cite{gao2026hgnnv2}, which is the reference point for any claim about its limitations.
No theoretical analysis of attention scoring in hypergraph operators has appeared since 2024;
the nearest treatments predate that, in work on equivariant hypergraph networks and on
generalising set transformers to hypergraphs \cite{kim2022equivariant, kim2021transformers}. The
transfer question raised in Section~\ref{sec:background:learning:attention} therefore remains
open on the theoretical side as well as the empirical one.

\section{Detection in Networked and Cyber-Physical Systems}
\label{sec:background:networks}

This section covers the unit of observation used in network detection, the families of method
defined over it, the particular character of industrial and IoT deployments, and the corpora on
which results are reported. Chapter~\ref{chap:anomaly} draws on it, although the Android
material precedes it for the reason given at the head of this chapter.

\subsection{Network Flows as a Unit of Observation}
\label{sec:background:networks:flows}

Detection at the flow level summarises the packets sharing a five-tuple into counts, durations
and byte totals. Working at that level lowers processing cost and survives payload encryption,
at the price of discarding content \cite{sperotto2010overview}. Network intrusion detection is
conventionally divided into signature matching against descriptions of known attacks and
anomaly detection against a learned model of normal behaviour \cite{paxson1999bro}, and the
second is what the rest of this section concerns. Section~\ref{sec:hybroid:7} states the
particular flow feature set Chapter~\ref{chap:hybroid} uses, rather than drawing it from here.

\subsection{Anomaly-Detection Method Families}
\label{sec:background:networks:families}

Anomalies are conventionally classed as point, contextual or collective, and detectors differ in
whether labels are available for both classes, for one, or for neither
\cite{chandola2009anomaly}. Classical methods score a point by the local density around it
\cite{breunig2000lof} or by how few partitions are needed to isolate it
\cite{liu2008isolation}. Deep detectors replace a hand-built feature space with a learned one,
and the shallow and deep families admit a common formulation \cite{ruff2021unifying,
pang2021deep}. Where the data are relational, the anomaly is a property of structure rather than
of any single record, and the object flagged may be a vertex, an edge or a subgraph
\cite{akoglu2015graph}. That observation is what connects this section to
Section~\ref{sec:background:graphs}.

The security setting resists the standard formulation. The cost of a false alarm is high,
network traffic is more variable than a model of normal behaviour admits, and data adequate for
evaluation are scarce \cite{sommer2010outside}.

\subsection{Communication Structure in Industrial and IoT Networks}
\label{sec:background:networks:industrial}

Industrial and IoT traffic is comparatively regular, because the devices are few and their
interactions are prescribed by the process under control. A whitelist of permitted communication
relations is therefore a workable model of normal behaviour. Barbosa et al.\ built such
whitelists in SCADA networks from four packet properties and reported both that the whitelists
stabilise and that rare legitimate flows continue to raise false positives
\cite{barbosa2013flow}; Yoon and Ciocarlie modelled communication patterns per host pair for the
same purpose \cite{yoon2014communication}. Neither of the two ideas that follow from this
originates with service-oriented deployments. Scoring behaviour against a graph of which host
talks to which was established for carrier traffic \cite{iliofotou2007network,
jin2009unveiling}, and the dependency graph between services was already being recovered
automatically from passively observed traffic in enterprise networks \cite{bahl2007towards,
chen2008automating, natarajan2012nsdminer}.

Where the system is cyber-physical, the observation may be the process rather than the traffic,
and detectors are then fitted to sensor and actuator series recorded from testbeds
\cite{mathur2016swat, kravchik2018detecting}. In IoT deployments the unit modelled is often the
individual device. Models of that kind hold under the conditions they were fitted to and lose a
substantial part of their measured performance when the deployment changes: reproduced figures
fall by roughly a quarter to a half beyond three months, and by a comparable margin between test
sites \cite{maali2025evaluating}.

\subsection{Datasets and Measurement Limits}
\label{sec:background:networks:datasets}

The corpora this field reports on are known to be defective. KDD~99 was shown to contain
redundant records that bias measured accuracy \cite{tavallaee2009nslkdd}, and a later audit of
CICIDS2017, UNSW-NB15 and comparable sets catalogues construction faults that allow trivial
features to separate attack traffic from benign \cite{flood2024bad}. Two further limits bear on
any figure reported here. Detectors reporting detection rates above 99 per cent have been shown
to be trained and tested on the same attack type, with performance not surviving attack types
absent from training \cite{kus2022false}. And when 70 published approaches were re-run on one
evaluation platform, four deliberately simple tests on process values, using minima, maxima and
rate of change, performed on par with them \cite{wolsing2022simple}.
Section~\ref{sec:background:evaluation} treats the protocol-level causes.

\section{Adversarial Machine Learning}
\label{sec:background:adversarial}

Chapter~\ref{chap:introduction} has already set out what the record shows about published
defences. This section supplies the vocabulary that account assumed, and it is the vocabulary
Chapters~\ref{chap:hybroid}, \ref{chap:hgann} and~\ref{chap:advml} use to state what they claim
and what they do not.

\paragraph{Threat models.} A threat model states four things: the attacker's goal, the
attacker's knowledge of the system, the attacker's capability to act on it, and the constraints
on what may be modified \cite{biggio2014security, vassilev2025adversarial}. A claim about an
attack or about robustness cannot be interpreted without all four \cite{carlini2019evaluating}.
Knowledge is conventionally divided into three levels. Under \emph{white box} the attacker knows
the training data, the feature representation, the learning algorithm, the trained parameters
and any deployed defence \cite{biggio2018wildpatterns}. Under \emph{black box} the attacker
knows neither parameters nor architecture and interacts through the prediction interface, either
by transferring from a substitute model trained on its outputs or by searching directly with
queries \cite{papernot2017practical}. \emph{Grey box} lies between the two: some components of
the white-box list are known and training cannot be influenced. Because the term is used
inconsistently in the literature, every use in this dissertation names which components the
attacker is assumed to know. One limit on the taxonomy should be recorded with it. A survey of
271 practitioners found that the knowledge assumptions of academic threat models are more
generous to the attacker than deployment conditions usually allow \cite{grosse2024towards}.

\paragraph{Classes of attack.} \emph{Evasion} modifies an input at test time so that a trained
model misclassifies it, leaving the model unchanged \cite{biggio2013evasion}. \emph{Poisoning}
manipulates the training data or the training process so that the resulting model is degraded or
carries behaviour the attacker chose \cite{biggio2012poisoning, jagielski2018manipulating}, and
the \emph{backdoor} is its sub-case, in which the implanted behaviour is conditioned on a trigger
\cite{gu2019badnets}. Confidentiality attacks form a third class: a functional copy of a
deployed model can be recovered through its prediction interface \cite{tramer2016stealing}. The
word adversarial covers all three and is not a synonym for evasion.

\paragraph{Feature space and problem space.} A \emph{feature-space} attack perturbs the feature
vector directly, with no requirement that any real input maps to the result. A
\emph{problem-space} attack modifies the artefact itself and must satisfy four constraints: the
transformation has to be available, semantics have to be preserved, the result has to be
plausible, and it has to survive preprocessing. Moving between the two is the inverse
feature-mapping problem, and in the malware setting the modified sample must still install and
still run \cite{pierazzi2020intriguing}. The distinction is load-bearing for
Chapters~\ref{chap:hybroid} and~\ref{chap:hgann}, because attacks on static Android detectors
operate at code level, upstream of feature extraction, so a representation derived from that
extraction inherits the attacker's influence whatever the learning architecture
\cite{song2026fcghunter}. Which operators reach which representation depends on what is
extracted, and each chapter settles that where it states its threat model. Early work on
Android established the feature-space case \cite{grosse2017adversarial}.

\paragraph{Adaptive attacks and robustness.} An \emph{adaptive} attack is constructed in
knowledge of the defence and tuned against that specific defence, in contrast to a fixed attack
applied unchanged \cite{tramer2020adaptive}. The recurring reason a defence appears to hold and
does not is gradient masking, where the defence obstructs the optimisation an attack relies on
without reducing the attacker's reach \cite{athalye2018obfuscated}. \emph{Robustness} is
therefore defined relative to a stated threat model and a stated attack budget. It is a property
of the pair, never of a model alone, and a robustness figure reported without those parameters
cannot be compared with another \cite{carlini2019evaluating}. Gradient-based attacks at fixed
perturbation budgets \cite{goodfellow2015explaining, madry2018pgd} are the reference instruments
against which the defences of Chapter~\ref{chap:advml} are measured.

\section{Performance Measures and Evaluation Methodology}
\label{sec:background:evaluation}

This section defines the measures used in Chapters~\ref{chap:anomaly} to~\ref{chap:advml}, with
explicit observation units and aggregation rules. Dimensionless scores are proportions:
$m$ corresponds to $100m$ per cent, and $m_1-m_2$ to $100(m_1-m_2)$ percentage points.

\subsection{Observation Units and Confusion Counts}
\label{sec:background:metrics:counts}

Let $\mathcal D=\{(x_i,y_i)\}_{i=1}^{N}$ be a labelled evaluation set with $N>0$, and let
$\widehat y_i=h(x_i)$ be the prediction of a fitted classifier. For a single-label task with
class set $\mathcal Y$, the confusion matrix $\mathbf C$ records true labels in its rows and
predicted labels in its columns:
\begin{equation}
  C_{ab}=\sum_{i=1}^{N}\mathbf1\{y_i=a,\ \widehat y_i=b\},
  \qquad a,b\in\mathcal Y.
  \label{eq:background:confusion}
\end{equation}
In binary detection, $1$ denotes an attack or malware and $0$ denotes normal or benign behaviour. With both axes
ordered as $(0,1)$,
\begin{equation}
  \mathbf C_{\mathrm{bin}}=
  \begin{pmatrix}\mathrm{TN}&\mathrm{FP}\\\mathrm{FN}&\mathrm{TP}\end{pmatrix},
  \qquad N=\mathrm{TP}+\mathrm{TN}+\mathrm{FP}+\mathrm{FN}.
  \label{eq:background:binary-confusion}
\end{equation}
True positives and true negatives are correct decisions. A false positive flags a negative
observation; a false negative misses a positive observation \cite{fawcett2006roc}.

The S7-300 evaluation counts network records, SWaT counts timestamps, and the Android studies
count applications. Timestamp metrics give longer attacks more weight; episode-level assessment
requires a separate matching rule and denominator \cite{sorbo2024navigating}. Collapsing attack
types into one positive class turns errors between those types into correct binary detections
(Section~\ref{sec:anomaly:iuno-results}).

\subsection{Binary Detection Measures and Class Prevalence}
\label{sec:background:metrics:binary}

Accuracy $A$ is the fraction of all observations classified correctly. Precision $P$ is the
fraction of positive predictions that are correct, whereas recall $R$ is the fraction of
positive observations detected:
\begin{equation}
  A=\frac{\mathrm{TP}+\mathrm{TN}}{N},\qquad
  P=\frac{\mathrm{TP}}{\mathrm{TP}+\mathrm{FP}},\qquad
  R=\frac{\mathrm{TP}}{\mathrm{TP}+\mathrm{FN}}.
  \label{eq:background:accuracy-precision-recall}
\end{equation}
Recall is also called sensitivity or the true-positive rate (TPR). The $F_1$ score combines
precision and recall through their harmonic mean:
\begin{equation}
  F_1=\frac{2\mathrm{TP}}{2\mathrm{TP}+\mathrm{FP}+\mathrm{FN}},\qquad
  H(P,R)=\frac{2PR}{P+R}=F_1
  \quad\text{when the ratios are defined}.
  \label{eq:background:f1}
\end{equation}
The count form omits true negatives and gives $F_1=0$ when $\mathrm{TP}=0$ and
$\mathrm{FP}+\mathrm{FN}>0$. We use the continuous extension $H(0,0)=0$ when both defined rates
are zero
\cite{fawcett2006roc,sorbo2024navigating}.

False-positive and false-negative rates (FPR and FNR) measure class-conditional errors.
Specificity, or true-negative rate (TNR), complements FPR; balanced accuracy averages the two class recalls
\cite{fawcett2006roc,grandini2020metrics}:
\begin{equation}
  \begin{aligned}
    \mathrm{FPR}&=\frac{\mathrm{FP}}{\mathrm{FP}+\mathrm{TN}},
    &\mathrm{TNR}&=\frac{\mathrm{TN}}{\mathrm{TN}+\mathrm{FP}}=1-\mathrm{FPR},\\
    \mathrm{FNR}&=\frac{\mathrm{FN}}{\mathrm{TP}+\mathrm{FN}}=1-R,
    &\mathrm{BAcc}&=\frac{R+\mathrm{TNR}}{2}.
  \end{aligned}
  \label{eq:background:error-rates}
\end{equation}
The false fraction among positive predictions is $1-P=\mathrm{FP}/(\mathrm{TP}+\mathrm{FP})$,
with a different denominator from FPR. Zero denominators make ratios undefined unless an
explicit replacement or exclusion rule is given.

Let $\pi=(\mathrm{TP}+\mathrm{FN})/N$ be the positive-class prevalence. Writing precision and
accuracy in terms of the class-conditional rates gives
\begin{equation}
  \begin{aligned}
    P(\pi)&=\frac{\pi\,\mathrm{TPR}}
      {\pi\,\mathrm{TPR}+(1-\pi)\,\mathrm{FPR}},\\
    A(\pi)&=\pi\,\mathrm{TPR}+(1-\pi)\,\mathrm{TNR}.
  \end{aligned}
  \label{eq:background:prevalence}
\end{equation}
Bayes' rule gives the precision identity \cite{axelsson2000base}. Varying $\pi$ at fixed TPR
and FPR gives a conditional prevalence-sensitivity calculation (Section~\ref{sec:hgann:10}).
Balanced accuracy removes the explicit prevalence weights, but changes in class-conditional
data distributions can still change both recalls.

\subsection{Multiclass and Multilabel Aggregation}
\label{sec:background:metrics:aggregation}

For $K$ mutually exclusive classes indexed by $1,\ldots,K$, let $n_k=\sum_b C_{kb}$ be the support of class $k$ and
$\pi_k=n_k/N$. A one-versus-rest calculation treats class $k$ as positive and all other classes
as negative:
\begin{equation}
  \begin{aligned}
    \mathrm{TP}_k&=C_{kk},
    &\mathrm{FN}_k&=n_k-C_{kk},\\
    \mathrm{FP}_k&=\sum_a C_{ak}-C_{kk},
    &\mathrm{TN}_k&=N-\mathrm{TP}_k-\mathrm{FP}_k-\mathrm{FN}_k.
  \end{aligned}
  \label{eq:background:one-vs-rest}
\end{equation}
The binary definitions then give $P_k$, $R_k$ and $F_{1,k}$. For any of these measures $M_k$,
macro averaging assigns equal class weights, while support-weighted averaging assigns weights
proportional to class size:
\begin{equation}
  M_{\mathrm{macro}}=\frac1K\sum_{k=1}^{K}M_k,
  \qquad
  M_{\mathrm{weighted}}=\sum_{k=1}^{K}\pi_k M_k.
  \label{eq:background:class-averages}
\end{equation}
Multiclass balanced accuracy is $R_{\mathrm{macro}}$. S7-300 results report both macro and
weighted $F_1$ to account for the dominant normal class. A high macro average does not establish
uniform per-class performance.

Micro averaging instead pools the one-versus-rest counts before calculating a ratio
\cite{grandini2020metrics}:
\begin{equation}
  P_{\mathrm{micro}}=
  \frac{\sum_k\mathrm{TP}_k}{\sum_k(\mathrm{TP}_k+\mathrm{FP}_k)},\qquad
  R_{\mathrm{micro}}=
  \frac{\sum_k\mathrm{TP}_k}{\sum_k(\mathrm{TP}_k+\mathrm{FN}_k)}.
  \label{eq:background:micro-averages}
\end{equation}
Micro $F_1$ uses the same count formula as Equation~\ref{eq:background:f1}, after summing over
classes. With one true and one predicted label per observation and all classes included,
both denominators in Equation~\ref{eq:background:micro-averages} equal $N$. With defined class
recalls, or an explicit zero-fill convention for classes of zero support, this gives
\begin{equation}
  P_{\mathrm{micro}}=R_{\mathrm{micro}}=F_{1,\mathrm{micro}}
  =R_{\mathrm{weighted}}=A=\frac{\operatorname{tr}(\mathbf C)}{N}.
  \label{eq:background:micro-identity}
\end{equation}
Opitz and Burst distinguish the mean of class $F_1$ scores from the harmonic mean of macro
precision and recall
\cite{opitz2019macrof1}. For defined class ratios,
\begin{equation}
  F_{1,\mathrm{macro}}
  =\frac1K\sum_k H(P_k,R_k)
  \leq H(P_{\mathrm{macro}},R_{\mathrm{macro}})
  \leq\frac{P_{\mathrm{macro}}+R_{\mathrm{macro}}}{2}.
  \label{eq:background:macro-f1}
\end{equation}
Here, macro $F_1$ denotes the mean of class scores. Unspecified source conventions remain
unresolved in Chapters~\ref{chap:hybroid} and~\ref{chap:hgann}. Reporting must identify the class
set and any zero-denominator policy.

Always predicting class $k$ gives accuracy $\pi_k$, while $\max_k\pi_k$ is the empirical
majority-class reference. A deployable baseline must choose its label before test evaluation,
for example from training labels, and may not match the test majority.

Multilabel prediction assigns several binary labels to each observation. In the facial-attribute
task, $y_i,\widehat y_i\in\{0,1\}^{J}$ and the reported accuracy averages individual label
decisions:
\begin{equation}
  A_{\mathrm{attr}}=\frac1{NJ}\sum_{i=1}^{N}\sum_{j=1}^{J}
    \mathbf1\{\widehat y_{ij}=y_{ij}\}.
  \label{eq:background:attribute-accuracy}
\end{equation}
Chapter~\ref{chap:advml} uses $J=40$ \cite{sparta2022d76}. Each attribute decision contributes
separately; the score does not require all labels of an image to be correct.

\subsection{ROC Curves and Area under the Curve}
\label{sec:background:metrics:roc}

For scores $s_i$ whose larger values favour the positive class, thresholding gives
$\widehat y_i(\tau)=\mathbf1\{s_i\geq\tau\}$. Varying $\tau$ traces the receiver operating
characteristic (ROC) points $(\mathrm{FPR}(\tau),\mathrm{TPR}(\tau))$
\cite{fawcett2006roc}. With $r(u)$ the linearly interpolated TPR at FPR $u$, the empirical area
under the curve (AUC) is
\begin{equation}
  \begin{aligned}
    \mathrm{AUC}&=\int_0^1 r(u)\,\mathrm du\\
    &=\frac1{N_+N_-}\sum_{i:y_i=1}\sum_{j:y_j=0}
      \left[\mathbf1\{s_i>s_j\}+\frac12\mathbf1\{s_i=s_j\}\right],
  \end{aligned}
  \label{eq:background:auc}
\end{equation}
where $N_+,N_->0$ count positive and negative observations. AUC measures positive--negative pair ranking,
with half credit for ties. Perfect and independent random rankings give one and expected
$0.5$, respectively. One confusion matrix gives only an operating point; AUC specifies neither
a deployment threshold nor probability calibration.

Multiclass one-versus-rest curves compare each class with the remainder. Unweighted macro AUC
is $K^{-1}\sum_k\mathrm{AUC}_k$; other weighting or curve-averaging conventions must be stated
\cite{fawcett2006roc}. \hybroid{} reports per-class ROC AUC (Chapter~\ref{chap:hybroid}).

\subsection{Performance under Adversarial Input}
\label{sec:background:metrics:adversarial}

An adversarial performance measure requires a fixed threat model and valid perturbed inputs
\cite{carlini2019evaluating}. For an attack $a$ applied to matched, single-label observations,
let $A_{\mathrm{clean}}$ be the clean accuracy and define
\begin{equation}
  A_a=\frac1N\sum_{i=1}^{N}\mathbf1\{h(a(x_i))=y_i\},
  \qquad \Delta_a=A_{\mathrm{clean}}-A_a.
  \label{eq:background:accuracy-loss}
\end{equation}
The SAFAIR score in Chapter~\ref{chap:advml} uses this signed accuracy difference
\cite{sparta2022d76}. For the attribute task, Equation~\ref{eq:background:attribute-accuracy}
supplies both terms. Report clean and attacked accuracy alongside the loss: equal differences
can occur at different levels of classification performance.

Attack success measures attainment of a specified goal. One explicit convention restricts
evaluation to initially correct examples, $I_{\mathrm{corr}}=\{i:h(x_i)=y_i\}$. For target
labels $t_i$, let $I_t=\{i\in I_{\mathrm{corr}}:t_i\neq y_i\}$. Conditional untargeted and
targeted attack-success rates are then
\begin{equation}
  \begin{aligned}
    \mathrm{ASR}_{\mathrm{untargeted}}
      &=\frac1{|I_{\mathrm{corr}}|}\sum_{i\in I_{\mathrm{corr}}}
        \mathbf1\{h(a(x_i))\neq y_i\},\\
    \mathrm{ASR}_{\mathrm{targeted}}
      &=\frac1{|I_t|}\sum_{i\in I_t}\mathbf1\{h(a(x_i))=t_i\}.
  \end{aligned}
  \label{eq:background:attack-success}
\end{equation}
Report eligible-set sizes and any alternative eligibility rule; empty sets give undefined rates.
These goal-specific rates differ from $\Delta_a$. Chapter~\ref{chap:advml} develops attribute
criteria and per-example attack-suite aggregation as analytical measures, not recovered contest
scores. Finite attacks do not certify worst-case robustness
\cite{carlini2019evaluating}.

\subsection{Aggregation and Operational Reporting}
\label{sec:background:metrics:reporting}

Let $M_b$ be a metric from fold or run $b$. Its mean and sample standard deviation are
\begin{equation}
  \overline M=\frac1B\sum_{b=1}^{B}M_b,
  \qquad
  s_M=\sqrt{\frac1{B-1}\sum_{b=1}^{B}(M_b-\overline M)^2},\qquad B>1.
  \label{eq:background:run-aggregation}
\end{equation}
Fold averaging can differ from calculating $M(\sum_b\mathbf C^{(b)})$ on pooled counts,
especially for nonlinear measures such as $F_1$. Class and fold averaging must be specified
separately. \hybroid{} reports five-fold means without the values needed for $s_M$.
Confidence intervals additionally require a sampling model that accounts for dependence;
correlated timestamps are not independent repetitions \cite{sorbo2024navigating}.

Operational performance uses time and resource denominators. For an observation interval of
duration $T$, let $n_{\mathrm{in}}$, $n_{\mathrm{proc}}$ and $n_{\mathrm{drop}}$ count input,
processed and dropped frames. The extractor measurements in Chapter~\ref{chap:anomaly} use
\begin{equation}
  q_{\mathrm{in}}=\frac{n_{\mathrm{in}}}{T},\qquad
  q_{\mathrm{proc}}=\frac{n_{\mathrm{proc}}}{T},\qquad
  p_{\mathrm{drop}}=\frac{n_{\mathrm{drop}}}{n_{\mathrm{in}}}.
  \label{eq:background:throughput}
\end{equation}
Using seconds gives frame rates per second. Chapter~\ref{chap:anomaly} also reports mean
whole-device CPU utilisation and peak memory. Its 95th-percentile delay from window close to
record emission excludes acquisition time; service-graph latency measures inter-service round trips. Timing endpoints
and summary statistics must be stated before comparison.

Two further ratios describe operational scope:
\begin{equation}
  r_{\mathrm{FA}}=\frac{n_{\mathrm{false\ alarm\ episodes}}}{T_{\mathrm{benign}}},
  \qquad
  c_{\mathrm{proc}}=\frac{N_{\mathrm{analysed}}}{N_{\mathrm{source}}}.
  \label{eq:background:alarm-coverage}
\end{equation}
The episode rate requires a grouping rule and benign exposure in a stated time unit; record FPR
alone cannot determine it. Processing coverage gives the analysed fraction of source applications.
Classification scores exclude the remaining applications, as in \hybroid{}. Complexity bounds
in Chapter~\ref{chap:hgann} describe scaling under stated assumptions, rather than measured
runtime or memory use.

\subsection{Experimental Protocol and Bias}
\label{sec:background:evaluation:bias}

Metric definitions fix the quantities calculated from an evaluation. The following considerations
concern how the data and protocol determine which conclusions those quantities support.

\paragraph{Sampling bias in space and time.} \emph{Temporal bias} is a split that allows a model
to learn from samples postdating the ones it is tested on, which reports a figure no deployment
reproduces. \emph{Spatial bias} is a test set whose composition differs from the deployment,
which leaves precision uninterpretable. A sound evaluation constrains both: training data must
precede test data, and the ratio between classes in the test set must reflect the setting the
system is claimed for \cite{pendlebury2019tesseract}. Drift is the mechanism that makes a
randomly ordered split optimistic \cite{jordaney2017transcend}. Curation of the corpus decides
as much as the split does, and guidance published since those constraints were formalised
revises parts of them \cite{chow2025beyond, kan2024tesseract}.

\paragraph{Leakage.} Any transformation fitted before the split moves information from the test
set into training. Resampling, feature selection, normalisation and model selection all have
this property when applied to the whole corpus, and the resulting figure measures the procedure
rather than the model \cite{arp2022dos}. Duplicate and near-duplicate samples are a second
channel, since repackaging produces applications that differ only in the injected payload and
copies of one application may fall on both sides of the split. The measured cost is uneven:
limited for supervised classification, and substantially larger for unsupervised family
clustering \cite{zhao2021impact, irolla2018duplication}.

\paragraph{Event definitions and benchmark labels.} In time-series detection,
\emph{point adjustment} marks an entire anomalous window as detected once
any single point inside it is flagged, which inflates F1 to the point that randomly generated
scores outperform published detectors \cite{kim2022towards, sorbo2024navigating}. Some standard
benchmarks also contain mislabelled or trivially separable anomalies, so a high score on them is
weak evidence \cite{liu2024elephant}.

\paragraph{Baselines and adaptive evaluation.} A reported gain is not established until a simple
baseline has been run under the same protocol. Where that has been done, the gap has often
disappeared: re-implementing eight provenance-based intrusion detectors in one framework left a
simple neural network matching the state of the art on most of the datasets tested
\cite{bilot2025sometimes}. For a defence the corresponding rule concerns the attack. A
robustness claim is bounded by the strongest attack applied to it, so the defence has to be
evaluated against an attack adapted to it \cite{carlini2019evaluating}, and defences evaluated
only by their own authors have repeatedly failed under adapted attacks
\cite{athalye2018obfuscated, tramer2020adaptive}.

\paragraph{The standard adopted here.} Every empirical chapter of this dissertation states its
split policy and uses time-aware splits where the timestamps permit, reports the class ratio
alongside any aggregate score, reports no point-adjusted figure, places a trivial baseline beside
each headline result, and states preprocessing and random seeds \cite{daoudi2021lessons}. Results
carried over from the source publications cannot be retrofitted to all of this, and each chapter
says in its Limitations which parts were not met.

\section{Summary and Research Gaps}
\label{sec:background:summary}

Three threads run through the preceding sections. The first concerns what an analysis system is
able to observe. An Android package exposes a manifest and bytecode to static reading, exposes
executed behaviour and recorded traffic to dynamic instrumentation, and neither mode dominates
the other. Network deployments expose flows, and industrial and IoT systems additionally expose
process values and a comparatively regular pattern of who communicates with whom. The second
thread concerns how those observations are structured. Program behaviour is encoded as
control-flow and function-call graphs, communication as graphs over hosts or services, and both
may be lifted to hypergraphs when a relation holds among more than two elements. The operators
that learn over these structures are bounded in what they can distinguish, and the projection
from a hypergraph back to a graph discards information that has now been measured. The third
thread concerns measurement. Reported performance depends on the split, on the corpus, on the
baseline it is compared against, and, where a defence is claimed, on the strength of the attack
applied to it.

Four questions are left open by this material, and each is taken up by one of the research
questions stated in Section~\ref{sec:intro:rqs}.

\phantomsection
\paragraph{G1. How behaviour in networked and industrial systems should be represented as
deployments change.}\label{gap:1}
Communication whitelists, service-dependency graphs, flow records, and process time series have
each supported anomaly detection under the conditions in which they were fitted
\cite{barbosa2013flow,yoon2014communication,goh2016swatdataset}. The evidence is split across
different observation units and system assumptions, and traffic-based models lose a substantial
part of their measured performance once a deployment moves in time or place
\cite{maali2025evaluating}. The literature lacks a common account of how these representation
choices support detection as the observed system changes. RQ1 takes up the question through
communication graphs and modular analysis of network and process features.

\phantomsection
\paragraph{G2. The value of combining program structure with observed behaviour.}\label{gap:2}
MARVIN and the study by Onwuzurike et al.\ report benefits from combining static and dynamic
features in within-system comparisons \cite{lindorfer2015marvin, onwuzurike2018family}.
Those results concern particular observation types and prediction tasks. Whether the benefit
persists when either changes requires comparison of the fused representation with each
modality alone on the same samples. RQ2 investigates this question for program-graph and
network-flow features in binary detection and malware-category classification.

\phantomsection
\paragraph{G3. Whether higher-order representation of program structure pays for itself.}\label{gap:3}
Program structure is represented pairwise, and a relation holding among several functions at
once is flattened into a set of edges that cannot record it as one group
\cite{zhou2006learning}. Hypergraphs remove the restriction and the cost of the projection has
been quantified \cite{kang2024low, wang2024graphs}, yet the sceptical result also stands: an
operator whose vertex weights do not vary between the hyperedges a vertex belongs to is
equivalent to one acting on a weighted clique expansion \cite{agarwal2006higherorder,
chitra2019randomwalks}. Whether higher-order encoding, and the attention weighting often applied
with it, improves detection and classification of Android malware is open, and no theoretical
analysis of attention scoring in hypergraph operators has appeared since 2024
\cite{kim2022equivariant, kim2021transformers}. RQ3 takes up this question.

\phantomsection
\paragraph{G4. How an adversary model becomes a reproducible benchmark.}\label{gap:4} Defences are commonly
evaluated against attacks chosen by the party proposing them, and a robustness claim is bounded
by the strongest attack applied \cite{carlini2019evaluating}. Threat models are stated
inconsistently, and the knowledge assumptions used in academic evaluations are more generous
than deployment usually allows \cite{grosse2024towards}. How an explicit adversary model is
translated into a benchmark that a third party can reproduce, and what such a benchmark then
shows about the relation between clean predictive performance and robustness, has no accepted
answer. RQ4 takes up this question.
